% GENERATED FILE. Edit canonical lessons, then run npm run generate:books.
% canonical-source-hash: fnv1a64:f8cb240e1fe1d179
% canonical-lessons: MW-C223-jhulno, MW-C223-hindolo, MW-C223-dhol, MW-C223-jhando, MW-C223-tamaso

\chapter{Off to the Fair}
\label{ch:mw-off-to-the-fair}
\hlchaptermodality{223}

\begin{chapteropening}
\textbf{By the end of this chapter:} \emph{I can say I will go to the fair, ride the swings and hear the drum.}

Complete the last lesson of chapter 223: I can say I will go to the fair, ride the swings and hear the drum.
\end{chapteropening}

\section[\texorpdfstring{\mw{झूलणो} (jhūlṇo)}{झूलणो (jhūlṇo)}]{\mw{झूलणो} (jhūlṇo) — to swing}

\label{lesson:MW-C223-jhulno}



\begin{quote}
Before the new one: say the Marwadi for to shiver, then the Marwadi for fog.
\end{quote}

\subsection*{You'll want to know: \mw{झूलणो}}
\textbf{\mw{झूलणो}} — \emph{jhūlṇo} — \textquotedblleft{}to swing\textquotedblright{}.

\begin{quote}
\mw{मेळो} \mw{काल} \mw{है} \mw{अर} \mw{हूं} \mw{झूलसूं।}
\end{quote}

\begin{quote}
\emph{meḷo kāl hai ar hū̃ jhūlsū̃.}
\end{quote}

\begin{quote}
The fair is tomorrow, and I'll ride the swings.
\end{quote}

Past (he …): \mw{वो} \mw{झूल्यो}, \emph{vo jhūlyo}.

\begin{cousinweb}
\textbf{\mw{झूलणो}} (\emph{jhūlṇo}) is inherited from a reconstructed \emph{*jhulyati}, \textquotedblleft{}swings\textquotedblright{}; CDIAL lists Marwari \emph{jhūlṇo}, \textquotedblleft{}to swing\textquotedblright{} (CDIAL 5406).
\end{cousinweb}

\begin{grammarlens}[title={What you've built}]
One more word for saying what comes next.
\end{grammarlens}

\subsection*{Your turn}
\begin{itemize}
\raggedright
  \item \emph{Say it:} \emph{jhūlṇo}
  \item \emph{Say it:} \emph{jhūlṇo}, once more
  \item \emph{Say it:} \emph{meḷo kāl hai ar hū̃ jhūlsū̃.} — The fair is tomorrow, and I'll ride the swings.
  \item \emph{Recall:} say the Marwadi for to shiver, then the Marwadi for fog, then say \emph{jhūlṇo} again
\end{itemize}

\begin{tcolorbox}[breakable,colback=teal!4,colframe=teal!35!black,title={Before you move on}]
What does \mw{झूलणो} mean? (\textquotedblleft{}To swing\textquotedblright{}.) Say it once more.
\end{tcolorbox}
\section[\texorpdfstring{\mw{हिंडोळो} (hiṇḍoḷo)}{हिंडोळो (hiṇḍoḷo)}]{\mw{हिंडोळो} (hiṇḍoḷo) — a big swing; a swing cradle}

\label{lesson:MW-C223-hindolo}



\begin{quote}
Before the new one: say the Marwadi for fog, then the Marwadi for to swing.
\end{quote}

\subsection*{You'll want to know: \mw{हिंडोळो}}
\textbf{\mw{हिंडोळो}} — \emph{hiṇḍoḷo} — \textquotedblleft{}a big swing; a swing cradle\textquotedblright{}.

\begin{quote}
\mw{हिंडोळो} \mw{घणो} \mw{बड़ो} \mw{है।}
\end{quote}

\begin{quote}
\emph{hiṇḍoḷo ghaṇo baṛo hai.}
\end{quote}

\begin{quote}
The swing is very big.
\end{quote}

A masculine noun.

\begin{cousinweb}
\textbf{\mw{हिंडोळो}} (\emph{hiṇḍoḷo}) is inherited from Sanskrit \emph{hindola}, \textquotedblleft{}a swing, a swing cradle\textquotedblright{}; CDIAL lists an Old Marwari word for a revolving swing under it (CDIAL 14094).
\end{cousinweb}

\begin{grammarlens}[title={What you've built}]
One more word for saying what comes next.
\end{grammarlens}

\subsection*{Your turn}
\begin{itemize}
\raggedright
  \item \emph{Say it:} \emph{hiṇḍoḷo}
  \item \emph{Say it:} \emph{hiṇḍoḷo}, once more
  \item \emph{Say it:} \emph{hiṇḍoḷo ghaṇo baṛo hai.} — The swing is very big.
  \item \emph{Recall:} say the Marwadi for fog, then the Marwadi for to swing, then say \emph{hiṇḍoḷo} again
\end{itemize}

\begin{tcolorbox}[breakable,colback=teal!4,colframe=teal!35!black,title={Before you move on}]
What does \mw{हिंडोळो} mean? (\textquotedblleft{}A big swing; a swing cradle\textquotedblright{}.) Say it once more.
\end{tcolorbox}
\section[\texorpdfstring{\mw{ढोल} (ḍhol)}{ढोल (ḍhol)}]{\mw{ढोल} (ḍhol) — a big drum}

\label{lesson:MW-C223-dhol}



\begin{quote}
Before the new one: say the Marwadi for to swing, then the Marwadi for a big swing.
\end{quote}

\subsection*{You'll want to know: \mw{ढोल}}
\textbf{\mw{ढोल}} — \emph{ḍhol} — \textquotedblleft{}a big drum\textquotedblright{}.

\begin{quote}
\mw{म्हैं} \mw{ढोल} \mw{सुण्यो।}
\end{quote}

\begin{quote}
\emph{mhaĩ ḍhol suṇyo.}
\end{quote}

\begin{quote}
I heard the drum.
\end{quote}

A masculine noun.

\begin{cousinweb}
\textbf{\mw{ढोल}} (\emph{ḍhol}) is inherited from Sanskrit \emph{ḍhōla}, \textquotedblleft{}a large drum\textquotedblright{} (CDIAL 5608, which lists Marwari \emph{ḍhol}).
\end{cousinweb}

\begin{grammarlens}[title={What you've built}]
One more word for saying what comes next.
\end{grammarlens}

\subsection*{Your turn}
\begin{itemize}
\raggedright
  \item \emph{Say it:} \emph{ḍhol}
  \item \emph{Say it:} \emph{ḍhol}, once more
  \item \emph{Say it:} \emph{mhaĩ ḍhol suṇyo.} — I heard the drum.
  \item \emph{Recall:} say the Marwadi for to swing, then the Marwadi for a big swing, then say \emph{ḍhol} again
\end{itemize}

\begin{tcolorbox}[breakable,colback=teal!4,colframe=teal!35!black,title={Before you move on}]
What does \mw{ढोल} mean? (\textquotedblleft{}A big drum\textquotedblright{}.) Say it once more.
\end{tcolorbox}
\section[\texorpdfstring{\mw{झंडो} (jhaṇḍo)}{झंडो (jhaṇḍo)}]{\mw{झंडो} (jhaṇḍo) — a flag (on a pole)}

\label{lesson:MW-C223-jhando}



\begin{quote}
Before the new one: say the Marwadi for a big swing, then the Marwadi for a big drum.
\end{quote}

\subsection*{You'll want to know: \mw{झंडो}}
\textbf{\mw{झंडो}} — \emph{jhaṇḍo} — \textquotedblleft{}a flag (on a pole)\textquotedblright{}.

\begin{quote}
\mw{मंदिर} \mw{माथै} \mw{लाल} \mw{झंडो} \mw{है।}
\end{quote}

\begin{quote}
\emph{mandir māthai lāl jhaṇḍo hai.}
\end{quote}

\begin{quote}
There's a red flag on the temple.
\end{quote}

A masculine noun.

\begin{cousinweb}
\textbf{\mw{झंडो}} (\emph{jhaṇḍo}) is inherited from a reconstructed \emph{*dhvajadaṇḍa}, \textquotedblleft{}a flagstaff\textquotedblright{}, built from Sanskrit \emph{dhvaja}, \textquotedblleft{}a flag\textquotedblright{}, and \emph{daṇḍa}, \textquotedblleft{}a staff\textquotedblright{}; CDIAL lists Hindi \emph{jhaṇḍā}, \textquotedblleft{}a flag\textquotedblright{}, and a Gujarati cousin. CDIAL gives no Marwari form in this entry (CDIAL 6898).
\end{cousinweb}

\begin{grammarlens}[title={What you've built}]
One more word for saying what comes next.
\end{grammarlens}

\subsection*{Your turn}
\begin{itemize}
\raggedright
  \item \emph{Say it:} \emph{jhaṇḍo}
  \item \emph{Say it:} \emph{jhaṇḍo}, once more
  \item \emph{Say it:} \emph{mandir māthai lāl jhaṇḍo hai.} — There's a red flag on the temple.
  \item \emph{Recall:} say the Marwadi for a big swing, then the Marwadi for a big drum, then say \emph{jhaṇḍo} again
\end{itemize}

\begin{tcolorbox}[breakable,colback=teal!4,colframe=teal!35!black,title={Before you move on}]
What does \mw{झंडो} mean? (\textquotedblleft{}A flag (on a pole)\textquotedblright{}.) Say it once more.
\end{tcolorbox}
\section[\texorpdfstring{\mw{तमासो} (tamāso)}{तमासो (tamāso)}]{\mw{तमासो} (tamāso) — a show; a spectacle}

\label{lesson:MW-C223-tamaso}



\begin{quote}
Before the new one: say the Marwadi for a big drum, then the Marwadi for a flag.
\end{quote}

\subsection*{You'll want to know: \mw{तमासो}}
\textbf{\mw{तमासो}} — \emph{tamāso} — \textquotedblleft{}a show; a spectacle\textquotedblright{}.

\begin{quote}
\mw{हूं} \mw{काल} \mw{तमासो} \mw{देखसूं।}
\end{quote}

\begin{quote}
\emph{hū̃ kāl tamāso dekhsū̃.}
\end{quote}

\begin{quote}
Tomorrow I'll watch the show.
\end{quote}

A masculine noun.

\begin{cousinweb}
\textbf{\mw{तमासो}} (\emph{tamāso}) is usually traced to Persian \emph{tamāšā}, \textquotedblleft{}a spectacle\textquotedblright{}, which Persian took from Arabic; Marwadi ends it in \emph{-o}.
\end{cousinweb}

\begin{grammarlens}[title={What you've built}]
One more word for saying what comes next.
\end{grammarlens}

\subsection*{Your turn}
\begin{itemize}
\raggedright
  \item \emph{Say it:} \emph{tamāso}
  \item \emph{Say it:} \emph{tamāso}, once more
  \item \emph{Say it:} \emph{hū̃ kāl tamāso dekhsū̃.} — Tomorrow I'll watch the show.
  \item \emph{Recall:} say the Marwadi for a big drum, then the Marwadi for a flag, then say \emph{tamāso} again
\end{itemize}

\begin{tcolorbox}[breakable,colback=teal!4,colframe=teal!35!black,title={Before you move on}]
What does \mw{तमासो} mean? (\textquotedblleft{}A show; a spectacle\textquotedblright{}.) Say it once more.
\end{tcolorbox}
